\subsection{补遗一：特征标层}

\label{ss:charsheaves}
\ssszh 本节开头略述了 dg 范畴的一些预备知识，随后讨论了范畴表示论中若干基本定义与构造，足以重新解释第 \ref{s:bbloc} 节预告过的一些结果。

本节余下部分，我们想再综述几条基本结果与构造。不过读者也可以直接跳到第 \ref{s:catloop} 节。

\ssszh 特征标层是群上（某些）伴随等变的 D-模。Lusztig 着眼于它们与有限李型群的不可约特征标之间的密切关联而开始了这项研究，尤其是在约化情形 \cite{luschar1, luschar23, luschar4, luschar5}。这是函数-层对应中一个格外丰富而不平凡的例子。

当然，有限群的特征标并不只是群上伴随等变的函数，而是从表示取迹得来的。人们可以问：特征标层是否有这样的解释？答案是有的，下面就来讨论。

\ssszh[对偶性] \label{ss:dualizable}回忆：要定义有限群表示的特征标，表示必须是有限维的，才能安全地取迹。在范畴情形中，「小到可以取迹」的类比物如下。

借助 $\DGCat$ 上的幺半结构，可以谈论可对偶化的 dg 范畴。在 $\DGCat$ 与 $\Vect$ 的类比下，它们对应于（有限维向量空间的有界复形）；确实，dg 范畴 $\sC$ 可对偶化当且仅当自然映射
%
\[
     \sC \rightarrow \udHom( \udHom(\sC, \Vect), \Vect)
\]
%
是等价。

对可对偶化的 dg 范畴，有典范等价 $\sC^\vee \simeq \udHom(\sC,
\Vect)$，即它在幺半意义下的对偶就是它的连续 dg-模。对任意 $\sS$，自然映射
%
\[
    \sS \otimes \sC^\vee \rightarrow \udHom(\sC, \sS)
\]
%
是等价。\footnote{把 $\sC^\vee$ 换成 $\Hom(\sC, \Vect)$，这又是可对偶化 dg 范畴的一条刻画。}更一般地，对任意 dg 范畴 $\sS$ 与 $\sS'$，有等价
%
\[
\Hom( \sS \otimes \sC, \sS') \simeq \Hom( \sS, \sC^\vee \otimes \sS').
\]
%

设 $\sA$ 是幺半 dg 范畴。由函子性，若 $\sC$ 是左 $\sA$-模，则 $\sC^\vee$ 是右 $\sA$-模。特别地，对可对偶化的 $\Dmod(K)$ 表示 $\sC$，它的对偶获得 $\Dmod(K)$ 的典范作用，即可以作反变表示。

在讨论如何定义可对偶化 $\Dmod(K)$ 表示的特征标层之前，我们先谈可对偶化范畴与模的一些基本例子。首先，下面这条一般事实保证了实践中遇到的多数 dg 范畴都可对偶化。

\begin{ex}若范畴 $\sC$ 紧生成，则它可对偶化。此外，它的对偶也紧生成，并且 $\sC$ 的紧对象与非余完备 dg 范畴 $\sC^\vee$ 的紧对象的反范畴之间有典范等价。
\end{ex}

上例有一些有用的特殊情形。

\begin{ex} 给定 $k$-代数 $A$，其左模范畴 $A\mod$ 可对偶化，对偶是右模范畴 $\on{mod-}A$。具体地说，求值映射
%
\[
   \rmod A \otimes A\mod \rightarrow \Vect
\]
%
把右模 $M$ 与左模 $N$ 送到它们的张量积 $M
\otimes_A N$。
\end{ex}

\begin{ex} \label{ex:dmodselfdual}给定簇 $X$，其 D-模范畴通过紧 D-模上的 Verdier 对偶而典范地自对偶。具体地说，所得的求值映射由复合
%
\[
 \Dmod(X) \otimes \Dmod(X) \simeq \Dmod(X \times X) \xrightarrow{\Delta^!}
 \Dmod(X) \xrightarrow{\pi_*} \Vect
\]
%
给出，其中与第 \ref{s:funselfdual} 小节一样，用 $\Delta: X
\rightarrow X \times X$ 记对角映射，用 $\pi: X \rightarrow \on{pt}$ 记到点的投影。

若 $X$ 是 $K$-簇，这个配对典范地 $\Dmod(K)$-等变，即把 $\Dmod(X)$ 与其反变表示等同起来。把 $\QCoh(X)$ 看作 $\QCoh(K)$ 的表示，作必要的调整后也有同样的陈述。
\end{ex}

\ssszh[特征标层] 现在可以说明如何取一个表示的特征标层。$\Dmod(K)$ 在范畴 $\sC$ 上的作用，等同于一个幺半函子
%
\[
\Dmod(K) \xrightarrow{\rho} \udHom(\sC, \sC).
\]
%
若 $\sC$ 可对偶化，则由等同 $\udHom(\sC, \sC) \simeq
\sC^\vee \otimes \sC$，它的自同态也可对偶化。于是可以把作用映射对偶化，得到矩阵系数的 $\Dmod(K \times K)$-等变函子
%
\[
   \udHom(\sC, \sC) \simeq \udHom(\sC, \sC)^\vee \xrightarrow{\rho^\vee}
   \Dmod(K)^\vee \simeq \Dmod(K).
\]
%

特别地，过渡到对角的 $K$-不变量，得到函子
%
\begin{equation} \label{e:mcequivfunc}
   \Hom_{\Dmod(K)}(\sC, \sC) \rightarrow \Dmod(K / K ).
\end{equation}
%
把它用于恒等自同态 $\id_{\sC}$，就得到 $K$ 上的伴随等变 D-模 $\chi_\sC$，即 $\sC$ 的\emph{特征标层}。

\ssszh 下面给出 $\chi_\sC$ 的更可动手的描述。给定可对偶化的范畴 $\sC$ 与自函子 $f: \sC \rightarrow \sC$，可以取它的迹，那是一个向量空间复形。与线性代数中完全一样，它由复合
%
\[
\Vect \xrightarrow{\id_{\sC}} \sC \otimes \sC^\vee \xrightarrow{f \otimes
\on{id}_{\sC^\vee}} \sC \otimes \sC^\vee \xrightarrow{\on{ev}} \Vect
\]
%
计算。也就是说，$\Vect$ 的这个复合自函子由与某个向量空间作张量给出，而这个向量空间按定义就是 $\on{tr}(f, \sC)$。要说明的是，当 $f$ 就是 $\sC$ 的恒等自同态时，所得的复形称作 $\sC$ 的 Hochschild 同调。

理解了在当下范畴情形中如何取迹之后，我们就能相当显式地描述特征标层了。也就是说，它的茎承载 $\sC$ 相应自同构的迹，与通常表示论中完全一样。即：对任意闭点 $g: \pt \rightarrow K$，有典范同构
%
\[
     g^!(\chi_\sC) \simeq \on{tr}( \delta_g \star -, \sC).
\]
%

\ssszh 下面提几个特点与例子。

\begin{ex} 可对偶化模之间的 $\Dmod(K)$-等变函子
$\sC \rightarrow \sD$
不一定诱导特征标层之间的映射。但是，若它带有连续的右伴随，则确实得到适当函子性的映射 $\chi_\sC
\rightarrow
\chi_{\sD}$。

要说明的是，在实际情形中 $\sC$ 是紧生成的，而「有右伴随」这一条件等价于把紧对象送到紧对象，也就是没有「小」的东西变成「大」的。
\end{ex}

\begin{ex} 范畴表示上的诱导函子与限制函子都保持可对偶化性质。在特征标上，它们重现特征标层通常的诱导函子与限制函子。
\end{ex}

\begin{ex} 若 $X$ 是 $K$-簇，它的特征标层可以如下描述。考虑 $K \times X$ 的闭子概形 $\on{Stab}$，它在 $X$ 的一点上的纤维就是该点的稳定化子。也就是说，$\on{Stab}$ 是作用映射与投影映射到 $X$ 的等化子。于是在平凡 $G$-等变映射
%
\[
\on{Stab} \rightarrow G
\]
%
下，对偶化层的推前就是 $\chi_{\Dmod(X)}$。特别地，由真基变换，$K$ 的任一元素在 $\Dmod(X)$ 上的迹就是其不动点的 Borel--Moore 同调。读者不妨把它与有限群、有限集上的类似断言作比较：在那里，群元素的迹给出不动点的个数。
\end{ex}

\begin{ex} 若 $K=G$ 是约化群，$\sC = \Dmod(\sB)$，则由上例，其特征标层是 $G/G$ 上的 Grothendieck--Springer 层。更一般地，等变自同态的矩阵系数函子（参见公式 \eqref{e:mcequivfunc}）形如
%
\[  \Dmod(B \backslash G / B) \rightarrow \Dmod(G/G).\]
%
这就重现了 Harish-Chandra 变换，或称极限圆变换，即沿对应 $$G/G \leftarrow G/B \rightarrow B \backslash G / B,$$ 的推-拉，其中左边与中间两项是关于伴随作用的商。后者正是 Lusztig 研究特征标层的出发点。
\end{ex}

下面有两条相关的注。

\begin{re}在当前的讨论中，我们把特征标层当作一个直接的构造：输入可对偶化的表示，输出伴随等变的 D-模。

另一种视角是把 $\Dmod(K/K)$ 解释为 $\Dmod(K)$ 的水平迹。确实，对任意幺半 dg 范畴 $\sA$，后者定义为 $\sA \otimes_{\sA \otimes \sA^{\on{op}}} \sA$；或者等价地，定义为 $\sA\mod$ 作为可对偶化 $\DGCat$-线性范畴的 2-范畴迹。若把范畴层级降低一层，即用一个代数 $A$ 代替幺半 dg 范畴，则这段讨论的类比物是把 $A\mod$ 的 Hochschild 同调实现为
%
\[
  A \otimes_{A \otimes A^{\on{op}}} A.
\]
%

在当下情形，利用范畴余不变量与商叠上 D-模之间的对应，有平凡的等同
%
\[
   \Dmod(K) \otimes_{\Dmod(K \times K)} \Dmod(K) \simeq \Dmod(K
   \overset{K \times K} \times K) \simeq \Dmod(K / K).
\]
%
一般的形式体系把可对偶化的 $\sA$-模对应到其水平迹的一个对象，而这就把上面关于特征标层的讨论作为特殊情形包含进来；进一步的讨论参见 \cite{gkrvtoyshtuka} 与 \cite{bznadlerchartheory}。
\end{re}

\begin{re}本节所讨论的理论还有平行的版本：把 $K$ 上的 D-模换成解析的或 \'etale 构造层，参见下文的第 \ref{ss:variants} 小节。当取定义在有限域 $\Fq$ 上的 $K$ 的 \'etale 层时，函数-层对应的一个版本把 $K(\Fq)$ 的模范畴等同于 $\on{Shv}(K)\mod$ 的 Frobenius 迹。这把范畴表示送到通常的表示，并且相应地，把特征标层送到通常的特征标；参见 \cite{Gaitsgoryinventshtukas}。
\end{re}
